\section{Conclusão e Próximas Etapas}
\label{sec:conclusao}

Este estudo demonstrou a viabilidade técnica e a pertinência jurídica de modelar o ciclo de habilitação, co-disputa licitatória, propriedade societária e execução orçamentária do município de Divinópolis/MG (exercício de 2025) por meio de um multigrafo heterogêneo multicamada dedicado (\texttt{01\_\allowbreak analise\_\allowbreak socios}). A partir de oito tabelas públicas do SICOM/TCE-MG integradas com 100\% de cobertura QSA nos 673 CNPJs válidos, o \textit{pipeline} reduziu 1.788 vértices e 5.914 arestas ($Q=0{,}7673$) a um subgrafo auditável de 26 vértices e 23 arestas, particionado em seis componentes independentes e três triângulos fechados $K_3$. O cruzamento financeiro pelos R\$\,140,29 milhões empenhados líquidos, a decomposição setorial e por item (padrão compatível com \textit{cover bidding} em \#4 e \#5) e o filtro anti-falso positivo separaram três núcleos críticos com evidência dura e contratos ativos (Componentes \#1, \#4 e \#2) de um alerta sem empenho (\#5), do estrato \texttt{MEDIO} podado (R\$\,9,11 milhões em 61 empresas distintas) e de dois casos puramente onomásticos (incluindo o Componente \#3, de R\$\,1,36 milhão), com execução integral em apenas 10,607\,s em memória.

Quatro limitações empíricas da base pública foram identificadas e tratadas ao longo da pesquisa:
\begin{enumerate}
  \item \textbf{Mascaramento de CPFs na origem:} o SICOM omite os dados nominais dos sócios e o QSA divulgado pela Receita Federal oculta os seis dígitos centrais do CPF (\texttt{***XXXXXX**}) e suprime o CPF de titulares MEI, além de registrar 1 documento mascarado com 12 co-disputas e R\$\,4.620,00 empenhados, tornando o casamento nominal dependente de transliteração e suscetível a homônimos;
  \item \textbf{Chave sentinela \texttt{-1} intermódulo:} como contratos e empenhos no SICOM 2025 trazem \texttt{seq\_\allowbreak licitacao = -1} e \texttt{homolog\allowbreak Licitacao\allowbreak .csv} cobre apenas 88 vencedores no exercício, a materialidade integrada por CNPJ/CPF higienizado reflete o volume contratual global do fornecedor em 2025, demandando conferência documental específica do processo na auditoria humana;
  \item \textbf{Latência de \textit{rate limit} em execução fria:} embora o processamento em grafos e auditoria execute em 10,607\,s, a consulta inicial de 673 CNPJs na BrasilAPI demanda no mínimo 673\,s devido ao intervalo de 1,0\,s por requisição, tornando o \textit{cache} local indispensável;
  \item \textbf{Sensibilidade onomástica e conformidade com a LGPD:} coincidências isoladas de sobrenomes não autorizam presunção de vínculo econômico, exigindo rebaixamento automático para revisão manual, postura estritamente indiciária e pseudonimização integral na publicação científica, com reversão offline restrita em \texttt{mapeamento\_\allowbreak pseudonimos\_\allowbreak restrito\allowbreak .csv} (\texttt{.gitignore}).
\end{enumerate}

Como desdobramentos futuros específicos para o módulo societário, propõe-se: (i) estender a análise para séries multianuais do SICOM (2020--2025), capturando alterações contratuais imediatamente anteriores a certames; (ii) cruzar os vértices empresariais e societários com cadastros sancionatórios (CEIS/CNEP) e vínculos de agentes públicos; (iii) acoplar os núcleos societários validados a análises comportamentais de lances e concentração orçamentária em arquiteturas integradas ou modelos preditivos em grafos; e (iv) concluir a publicação do conjunto de dados pseudonimizado com atribuição de DOI definitivo no Zenodo~\cite{projeto_dataset_zenodo}, assegurando reprodutibilidade plena em ambiente Linux~\cite{projeto_github}.
